\begin{exo}{}
	On considère le triangle ci-dessous, dans lequel les côtés dépendent d'un nombre réel $x > 1$.

	\begin{center}
		\begin{tikzpicture}[scale=0.9, >=Stealth]
			\coordinate (A) at (0,0);
			\coordinate (B) at (6,0);
			\coordinate (C) at (4,3);
			\draw[very thick] (A) -- (B) node[midway, below=1mm] {$2x - 1$} -- (C) node[midway, above right=1mm] {$3x + 6$} -- cycle node[midway, above left=1mm] {$x + 9$};
			\node[left=1mm] at (A) {$A$};
			\node[right=1mm] at (B) {$B$};
			\node[above=1mm] at (C) {$C$};
		\end{tikzpicture}
	\end{center}
	\begin{enumerate}
		\item Pour quelle valeur de $x$ a-t-on $AB = AC$ ? Pour cette valeur de $x$, quelle est la longueur de chacun des côtés de $ABC$ ?
		\item Déterminer toutes les valeurs de $x$ pour lesquelles le triangle $ABC$ est isocèle.
		\item Peut-on trouver une valeur de $x$ pour laquelle le triangle $ABC$ est équilatéral ?
	\end{enumerate}
\end{exo}
